% Ad Atticum 5.15 --- headnote
% ad-atticum-05-15, 3 August 51 BC, written at Laodicea

\textit{Cicero to Atticus, written from Laodicea on the third day
before the Nones of August (3 August) 51 BC, two days after his
arrival at the capital of the conventus on 31 July. He notes the
date as the start of the year's count --- the year of imperium he
is determined will not be extended --- and immediately admits
what governing this provincial backwater feels like to a man
whose horses are still bred for Rome: the field is too small for
him, his ``splendid work goes to waste,'' he is judging cases at
Laodicea while A. Plotius judges them in the Forum, and he holds
the name only of two skeleton legions while Pompey (``our
friend'') has armies enormous in the comparison. The list of
what he misses --- daylight, Forum, City, his house, Atticus
himself --- is one of the most economical homesicknesses in the
correspondence.}

\textit{The second paragraph turns to expenses and to the running
joke of the trip. Cicero is being scrupulously abstinent by
Atticus's precepts --- so abstinent, he says, that the loan he
got from Atticus on exchange may have to be repaid by a fresh
loan; the wounds Appius Claudius left in the province are
visible everywhere, though Cicero is decent enough not to pull
them open. He is on the march from Laodicea into Lycaonia, on
toward Taurus, where he intends, with mock-military pomp, to
fight a pitched battle with one Moeragenes over a runaway slave
of Atticus's. The Latin slips into the iambic of an old
proverb, ``the pack-saddle on the ox is set'' --- a soldier's
load fastened on a beast born for the plough --- which is
exactly how the governorship feels. The closing instructions on
the postal route, through the masters of pasturage-dues and
harbour-dues of the publicans, show how the tax-farmers'
infrastructure carried the orator's letters home.}
